% ============================================================
%  LAS TRES PREGUNTAS DE PALMER, TRADUCIDAS A FILAS DEL MAPA DE DECISION
%  Fragmento aparte, para no mover la numeracion del resto del deck.
%  Se activa quitando el % de la linea \input que esta en semana01.tex,
%  justo despues del frame "Que preguntas motivaron todo este trabajo".
%  Sirve igual de repaso al abrir una sesion posterior: se da una pregunta
%  en lenguaje de biologo y el grupo dice que renglón del mapa le toca.
% ============================================================
\begin{frame}{Cómo se traduce una pregunta en una técnica}
  Ante cualquier pregunta que traiga alguien del área, un estadístico hace tres
  preguntas antes de mirar el mapa de decisión:

  \vspace{0.3cm}
  \begin{enumerate}
    \item ¿Cuál es la variable \textbf{respuesta}, y de qué \textbf{tipo} es?
      {\small (cuantitativa o categórica)}
      \vspace{0.15cm}
    \item ¿Cuál es la variable \textbf{explicativa}, y de qué tipo? Si es
      categórica, ¿\textbf{cuántos niveles} tiene?
      \vspace{0.15cm}
    \item ¿\textbf{Cómo se obtuvieron} los datos? ¿Grupos independientes, o dos
      medidas del \emph{mismo} individuo?
  \end{enumerate}

  \vspace{0.35cm}
  Las dos primeras eligen el renglón. La tercera decide entre dos renglones que se
  parecen. Y de fondo hay una cuarta: ¿la pregunta es de \textbf{diferencia}, de
  \textbf{relación} o de \textbf{predicción}?

  \vspace{0.3cm}
  \destaca{El mapa de decisión \textbf{no se consulta con las palabras del
  problema, sino con los tipos de variable}. Traducir de un lenguaje al otro es
  la habilidad que se practica todo el semestre.}
\end{frame}
% ============================================================
\begin{frame}{Las tres preguntas de Palmer, traducidas}
  \begin{center}
  {\scriptsize
  \begin{tabular}{@{}p{4.3cm}p{4.5cm}p{3.6cm}@{}}
    \toprule
    \textbf{Lo que dice quien trae el problema} & \textbf{Lo que oye un estadístico} &
    \textbf{Renglón del mapa} \\
    \midrule
    ¿Difieren machos y hembras en tamaño? &
      Cuantitativa contra categórica de \textbf{2 niveles}, grupos independientes &
      Dos grupos independientes {\scriptsize\color{grisT}(unidad 8)} \\
    \dots y si fueran los dos del \emph{mismo nido} &
      Cuantitativa, \textbf{dos medidas del mismo par} &
      Datos pareados {\scriptsize\color{grisT}(unidad 8)} \\
    ¿Y en cuál especie la diferencia es mayor? &
      Cuantitativa contra \textbf{dos} categóricas: es una \textbf{interacción} &
      \textcolor{rojoAc}{No está en el mapa: ANOVA de dos factores} \\
    ¿Se puede adivinar el sexo a partir de las medidas? &
      La respuesta ahora es \textbf{categórica}, con explicativas cuantitativas &
      \textcolor{rojoAc}{No está en el mapa: regresión logística} \\
    ¿Se relacionan el hielo marino y la alimentación? &
      Dos cuantitativas &
      Correlación {\scriptsize\color{grisT}(unidad 3)} y \Rc{lm()}
      {\scriptsize\color{grisT}(unidad 11)} \\
    \bottomrule
  \end{tabular}}
  \end{center}

  \vspace{0.3cm}
  \raggedright
  \destaca{\small Dos observaciones que valen todo el curso.

  \vspace{0.15cm}
  \textbf{El renglón 1 y el 2 son la misma pregunta con datos distintos.} El estudio
  midió al macho \emph{y} a la hembra del mismo nido. Eso es pareado, y la prueba
  pareada detecta diferencias más pequeñas. Pero nuestra tabla no trae el
  identificador del nido, así que perdimos esa ventaja.

  \vspace{0.15cm}
  \textbf{El renglón 1 y el 4 usan las mismas dos variables, con los papeles
  invertidos.} ``¿La masa depende del sexo?'' y ``¿puedo adivinar el sexo si conozco
  la masa?'' no son la misma pregunta, y no se contestan con la misma técnica.}
\end{frame}
